\subsection{Vogel's Approximation}
We can use vogels aproximation, it works in the following way, following uge 8 "Stepping Stone metoden og Vogel’s approksimation":

In Vogel’s Approximation Method, an edge is selected iteratively and then assigned as much flow as possible.

The edge is chosen as the least-cost edge in a row or column for which the difference between the cheapest and the second-cheapest cost in that row or column is maximal. This is done such that we save the maximal amount by choosing the cheapest edge to the fullest extend possible

The value (flow) assigned to the edge is the maximum possible amount, determined by the supplier’s remaining capacity and the customer’s remaining demand (i.e., the minimum of the two).

Since in each iteration either a remaining supply or a remaining demand is reduced to zero, there can be at most n+m−1 iterations in a balanced transportation problem.

According to the pseudocode, however, there are no further choices to make once only a single supplier or a single customer remains. This situation occurs after at most n+m−3 iterations.